\chapter{Клеточный автомат: как всё собрано}\label{ch:automaton}

\section{Поле, сетка, оси}

Модель живёт в поперечном сечении стенки: прямоугольник $240\times240$~мкм (по радиусу и по
окружности), покрытый клетками 0.4~мкм ($600\times600$). Поле периодическое по обеим осям ---
это условие Фурье. Ось $\theta$ (окружность, TD) идёт по столбцам, ось $r$ (радиус, ND) --- по
строкам. Окружное напряжение $\sigma_\theta$ действует вдоль столбцов.

\begin{warnbox}
Сечение двумерное: пластинка бесконечна вдоль оси трубы, наклон её плоскости из сечения не
учитывается, осевое напряжение входит только через среднюю $\eps_{zz}$. Отдельный трёхмерный
движок (\code{ca3d\_kinetic.py}) есть, но основная калибровка --- в 2D (гл.~\ref{ch:limits}).
\end{warnbox}

\section{Зёрна и текстура}

\fig{F19_Grains}{Зёрна на поле $30\times30$~мкм. В каждом --- след базисной плоскости, то есть
направление, вдоль которого ляжет пластинка (цвет --- крутизна). Оранжевым выделены грани,
на которых возможен межзёренный гидрид: их след близок к базисному следу одного из соседей.}

\begin{itemize}
\item \textbf{Зёрна} --- ячейки Вороного в «растянутой» метрике: семена разбросаны случайно, клетка
  принадлежит ближайшему семени, расстояние меряется с весами $1/a_r$ и $1/a_\theta$, так что зерно
  вытянуто вдоль окружности ($2.5\times4.5$~мкм у листа Lepine).
\item \textbf{Ориентация зерна} --- угол $\psi$ следа базисной плоскости к окружности. Ось $c$
  наклонена к радиусу на $\pm\chi_0$ ($\chi_0=30^\circ$ у Lepine, $38.5^\circ$ у Э635) с нормальным
  разбросом $\chi_s=26^\circ$; знак --- случайный.
\item \textbf{Грани} --- клетки на стыке двух ячеек Вороного; у каждой грани свой след. На ней
  разрешено межзёренное зарождение, если след отличается от базисного следа соседа не больше
  чем на $15^\circ$ (гл.~\ref{ch:nucleation}).
\end{itemize}

\begin{exercise}
Сколько зёрен $2.5\times4.5$~мкм в поле $240\times240$~мкм? Какая доля зёрен имеет крутой след
($|\psi|>65^\circ$) при $\chi_0=30^\circ$, $\chi_s=26^\circ$? (Подсказка: $|\psi|>65^\circ$ значит,
что ось $c$ отклонена от радиуса больше чем на $65^\circ$.)
\end{exercise}

\section{Пластинка в модели}

Пластинка --- отрезок с центром, углом $\psi$ (базисный след зерна или след грани), полудлиной
и толщиной $h=0.6$~мкм. На сетку она ложится долями: в каждой клетке --- доля площади пластинки,
посчитанная по $3\times3$ подточкам. Из долей складываются собственная деформация (с ореолом,
гл.~\ref{ch:plasticity}) и карта занятости; новая пластинка не может лечь на клетку, соседнюю
с занятой.

\section{История температуры}

\begin{enumerate}
\item \textbf{Нагрев} до $T_{\max}$: растворяется $\min(H,\TSSD(T_{\max}))$. Если растворилось не всё,
  остаток --- это самые длинные пластинки предварительно посчитанной окружной структуры (при
  нагреве первыми растворяются мелкие).
\item \textbf{Охлаждение} со скоростью $\dot T$ шагами не больше 0.5\,$^\circ$C (при всплеске
  зарождения --- мельче) до $T_{\text{кон}}=100\,^\circ$C, где почти весь водород уже выпал.
\end{enumerate}

\section{Шаг по времени}

\fig[0.95]{F20_Algorithm}{Один шаг автомата. Сначала --- поле и движущая сила по текущему
состоянию, затем случайные зародыши, рост, диффузия и остывание.}

Псевдокод шага:
\begin{codebox}\small
\textbf{для} каждой клетки $\mathbf x$:\\
\quad $g(\mathbf x)=\sigma(\mathbf x):\eps^*_{\text{нов}}(\mathbf x)/\eps_n$ \hfill (поле: пластинки + ореолы)\\
\quad $\Delta(\mathbf x)=\Delta_0+n_HRT\ln\big(c(\mathbf x)/c_{\TSSP}(T)\big)+\eps_n g(\mathbf x)+n_H\frac{Q}{T}\,\delta T(\mathbf x)$\\
\quad $r(\mathbf x)=\nu_c\exp\!\big[-B\big((\Delta_0/\Delta)^2-1\big)\big]$, если $\Delta>0$ и клетка свободна\\
$\lambda=\sum_{\mathbf x}r(\mathbf x)\,\delta t$;\ если $\lambda>6$ --- уменьшить $\delta t$\\
$N\sim\text{Пуассон}(\lambda)$;\ места --- с вероятностями $\propto r(\mathbf x)$\\
\textbf{для} каждого места: пластинка по базису зерна (или по грани) длиной $L_{\min}$,\\
\quad если в окрестности хватает водорода сверх TSSD; снять водород\\
решить упругость для новых пластинок (с ореолами) и прибавить к $\sigma$\\
вырастить кончики: $v=k_{\text{tip}}\,D/h\cdot(c-c_{\TSSD})/C_{\text{гидр}}$; продлить через границы\\
диффузия: $\hat c(\mathbf k)\leftarrow\hat c(\mathbf k)\,e^{-Dk^2\delta t}$\\
$T\leftarrow T-\dot T\,\delta t$
\end{codebox}

\begin{idea}
Ни одно правило автомата не говорит «клади пластинку по радиусу». Ориентацию задают зерно или
грань, а \emph{где} появится следующая пластинка и \emph{какого} семейства --- решает конкуренция
химии, поля соседей и нагрузки. Радиальные стопки и окружные пакеты --- следствие, а не вход.
\end{idea}

\begin{exercise}
Охлаждение с 415 до 100\,$^\circ$C со скоростью 3\,$^\circ$C/мин шагами 0.5\,$^\circ$C. Сколько шагов
и сколько минут модельного времени? Сколько это шагов при 0.3\,$^\circ$C/мин и как изменится
время счёта?
\end{exercise}

\section{Параметры по умолчанию}

\begin{table}[htbp]\centering\small
\caption{Основные параметры кинетического автомата (калибровка по Zircaloy-4, 178 ppm).}
\begin{tabular}{llll}\toprule
параметр & значение & смысл & откуда\\\midrule
поле, сетка & $240\times240$, 0.4~мкм & сечение $r$--$\theta$ & выбор\\
зерно & $2.5\times4.5$~мкм & размер по $r$ и $\theta$ & Lepine (EBSD)\\
$\chi_0$, $\chi_s$ & $30^\circ$, $26^\circ$ & наклон оси $c$ и разброс & Кернс, Lepine\\
$\eps_n$, $\eps_t$ & 0.072, 0.0458 & несоответствие & Carpenter 1973\\
$E$, $\nu$ & 90~ГПа, 0.34 & упругость & справочник\\
TSSD, TSSP & табл.~в гл.~\ref{ch:solubility} & линии растворимости & Плясов 2023\\
$B$, $\Delta_0$ & 80, 15~МПа & барьер и запас на TSSP & \textbf{подогнано}\\
app\_dT & $0.08\,^\circ$C/МПа & сдвиг выпадения нагрузкой & Vizcaíno 2014\\
bias\_dT & 10.9--12.4\,$^\circ$C & фора зёрен с осью $c$ по $r$ & \textbf{подогнано} (порог Cinbiz)\\
gb\_dT, gb\_tol & 1.5\,$^\circ$C, $15^\circ$ & фора и правило границы & теория, Qin/Son\\
$D_0$, $Q_D$ & $7.9\cdot10^{-7}$~м$^2$/с, 44.4~кДж/моль & диффузия & справочник\\
$h$, $L_{\min}$, $L_{\max}$ & 0.6, 1.2, 6~мкм & пластинка & допущения\\
$k_{\text{tip}}$, cross\_tol & 0.05, $15^\circ$ & рост, переход через границу & допущения\\
halo, free\_z & да / да & ореолы, свободная ось & новое (гл.~\ref{ch:plasticity})\\\bottomrule
\end{tabular}
\end{table}

\begin{codebox}
\textbf{Где в коде.} \code{ca\_kinetic.py} --- \code{KParams}, \code{run\_kinetic};
\code{ca\_hydride.py} --- \code{Params}, \code{make\_grains}, \code{plate\_eigen}, \code{Elastic};
\code{thermo.py} --- линии растворимости; \code{tool.py} --- «инструмент»: карточка материала, текстура
и история охлаждения $\to$ структура и меры. Запуск одного расчёта:
\code{run\_kinetic(KParams(sigma\_app=150, H\_ppm=178, ...))}; перебор --- \code{calib\_kin.py}.
Поле $240\times240$~мкм считается 2--7 минут на одном ядре.
\end{codebox}
